\chapter{Практичні поради для громадян}
\label{apdx:practical}

\subsection*{Як розпізнати шахрайство}

\begin{enumerate}[label=\arabic*)]
\item \parheading{Несподівані запити грошей} від знайомих — \keyword{передзвоніть іншим каналом}
\item \parheading{Помилки в доменах}: \example{pr1vat24.com} замість \example{privatbank.ua}
\item \parheading{Тиск і терміновість}: «блокують через 1 годину», «приз лише сьогодні»
\item \parheading{Запити паролів/PIN/CVV} — \important{легітимні організації ніколи не запитують}
\item \parheading{Висока дохідність} понад 30--50\% річних
\item \parheading{Орфографічні помилки} в офіційних листах
\item \parheading{Вимога оплати криптовалютою} або нерегульованими сервісами
\end{enumerate}

\subsection*{Перевірка веб-сайтів}

\parheading{Перед введенням даних:}
\begin{itemize}[label=--]
\item \important{HTTPS і замок} в адресному рядку обов'язкові
\item \term{Уважно читайте домен}: шахраї використовують схожі назви (typosquatting)
\item \keyword{Короткі URL розгортайте} через CheckShortURL
\end{itemize}

\parheading{Сервіси перевірки:}
\begin{itemize}[label=--]
\item \term{VirusTotal} — сканування посилань
\item \term{Google Safe Browsing} — база фішингу
\item \term{WHOIS} — власник домену
\end{itemize}

\subsection*{Захист фінансів}

\parheading{Банківські картки:}
\begin{itemize}[label=--]
\item \important{Увімкніть SMS-сповіщення} про всі операції
\item \term{Встановіть ліміти} на онлайн-платежі
\item \keyword{Віртуальні картки} для інтернет-покупок
\item \important{Блокуйте негайно} при підозрі: гаряча лінія 24/7
\end{itemize}

\parheading{Паролі та 2FA:}
\begin{itemize}[label=--]
\item \important{Мінімум 12 символів}, унікальні для кожного сервісу
\item \term{Менеджери паролів}: Bitwarden, 1Password, KeePassXC
\item \keyword{Увімкніть 2FA} на банках, поштi, соцмережах — краще застосунки (Google Authenticator), ніж SMS
\end{itemize}

\subsection*{Безпека пристроїв}

\parheading{Комп'ютер:}
\begin{enumerate}[label=\arabic*)]
\item Антивірус і регулярні оновлення баз
\item Оновлення ОС — \important{критичні патчі виправляють вразливості}
\item Брандмауер увімкнено
\item Блокувальники реклами (uBlock Origin)
\end{enumerate}

\parheading{Смартфон:}
\begin{itemize}[label=--]
\item \important{PIN/біометрія} для розблокування
\item \term{Віддалене стирання} (Find My Device)
\item \keyword{Лише офіційні магазини} — Google Play, App Store
\item VPN у публічних Wi-Fi
\end{itemize}

\subsection*{Безпека комунікацій}

\parheading{Електронна пошта:}
\begin{itemize}[label=--]
\item \important{Перевіряйте адресу}: \example{support@paypal.com} правильно, \example{support@paypa1.com} — фішинг
\item \keyword{Наводьте курсор на покликання} перед кліком
\item \term{Не відкривайте вкладення} \texttt{.exe}, \texttt{.scr}, \texttt{.bat}, \texttt{.js}
\end{itemize}

\parheading{Месенджери:}
\begin{itemize}[label=--]
\item \important{Обмежте доступ} до профілю (налаштування приватності)
\item \term{Перевіряйте запити від «знайомих»} дзвінком — акаунти зламують
\item \keyword{Не переходьте за підозрілими покликаннями}
\end{itemize}

\parheading{Телефонні дзвінки:}
\begin{itemize}[label=--]
\item \important{Банки ніколи не просять паролі чи CVV}
\item \term{Перевіряйте номер} на офіційному сайті
\item Додатки-визначники (Truecaller) фільтрують спам
\end{itemize}

\subsection*{Що робити, якщо стали жертвою}

\begin{enumerate}[label=\arabic*)]
\item \important{Зупиніть втрати}: заблокуйте картки, змініть паролі, вимкніть пристрій від інтернету

\item \important{Зафіксуйте докази}: скриншоти, номери/адреси шахраїв, дата й обставини

\item \important{Повідомте}:
   \begin{itemize}[label=--]
   \item \term{Кіберполіція}: \texttt{cyberpolice.gov.ua}, тел. 102
   \item Банк (для оскарження транзакції — протягом 24 год)
   \item Платформу (OLX, Facebook)
   \end{itemize}

\item \parheading{Заява до поліції} — статті 190, 361, 361-2 ККУ (див. Додаток Б)

\end{enumerate}

\subsection*{Корисні ресурси}

\parheading{Державні служби:}
\begin{itemize}[label=--]
\item \term{Кіберполіція}: \texttt{cyberpolice.gov.ua}, 0-800-505-170
\item \term{НБУ}: \texttt{bank.gov.ua}, 0-800-505-240
\item \term{CERT-UA}: \texttt{cert.gov.ua}
\end{itemize}

\parheading{Перевірка безпеки:}
\begin{itemize}[label=--]
\item \term{Have I Been Pwned}: перевірка на витік даних
\item \term{VirusTotal}: сканування файлів і посилань
\end{itemize}

\vspace{1em}

\noindent\important{Пам'ятайте}: \keyword{найкраща безпека — це критичне мислення}. Перевіряйте, не довіряйте сліпо, захищайте дані.
